\documentclass[12pt,a4paper]{article}
\usepackage[utf8]{inputenc}
\usepackage[spanish]{babel}
\usepackage{geometry}
\geometry{top=2.5cm, bottom=2.5cm, left=3cm, right=3cm}
\usepackage{graphicx}
\usepackage{hyperref}
\usepackage{xcolor}
\usepackage{fancyhdr}
\usepackage{amsmath}

\definecolor{unabblue}{HTML}{003865}
\definecolor{unaborange}{HTML}{E85C0B}
\definecolor{codegray}{rgb}{0.95,0.95,0.95}
\definecolor{codeorange}{rgb}{0.8,0.3,0.0}

\pagestyle{fancy}
\fancyhf{}
\fancyhead[L]{\textcolor{unabblue}{\textbf{Ciencia de Datos - UNAB}}}
\fancyhead[R]{\textcolor{unaborange}{\textbf{Taller Semana 3}}}
\fancyfoot[C]{\thepage}

\begin{document}

\begin{center}
    {\huge \textcolor{unabblue}{\textbf{Taller 3: Vectorización, Numpy y Exploración con Pandas}}} \\[0.5cm]
    {\large Docente: Dudbil Olvasada Pabon Riaño} \\[0.2cm]
    \rule{\textwidth}{1pt}
\end{center}

\vspace{0.5cm}

\section*{\textcolor{unabblue}{1. Objetivo del Taller}}
Demostrar la superioridad en rendimiento de la vectorización matemática frente a la iteración clásica de Python. El estudiante importará y explorará datos tubulares utilizando Pandas y ejecutará operaciones masivas a nivel de CPU usando Numpy.

\section*{\textcolor{unabblue}{2. Instrucciones Generales}}
Desarrolle los siguientes retos en un Jupyter Notebook (Google Colab). En este taller \textbf{SÍ} está permitido y es obligatorio el uso de \texttt{numpy} y \texttt{pandas}. Documente su proceso utilizando celdas de Markdown. Al finalizar, descargue el notebook o suba el enlace a su repositorio.

\vspace{0.5cm}
\noindent\fbox{%
    \parbox{\dimexpr\textwidth-2\fboxsep-2\fboxrule\relax}{%
        \textbf{\textcolor{unabblue}{Misión Analítica: IoT y Big Data}}\\[0.2cm]
        Una red de sensores IoT instalada en tractores agrícolas envía miles de mediciones de temperatura y humedad por minuto. El código actual en Python puro está colapsando el servidor por su lentitud. Su misión es refactorizar el procesamiento a Numpy y Pandas.
    }%
}
\vspace{0.5cm}

\subsection*{Reto 1: La Muerte del Bucle \texttt{for} (30\%)}
Para simular el volumen de datos, cree una lista en Python con 1 millón de registros aleatorios de temperatura (entre 10.0 y 40.0 grados Celsius):
\begin{center}
    \texttt{import random}\\
    \texttt{temp\_celsius = [random.uniform(10.0, 40.0) for \_ in range(1000000)]}
\end{center}
\begin{enumerate}
    \item \textbf{El método lento:} Escriba una función clásica con un bucle \texttt{for} que convierta todas las temperaturas de Celsius a Fahrenheit ($F = C \times 1.8 + 32$). Mida el tiempo de ejecución usando el módulo \texttt{time}.
    \item \textbf{La vía Numpy:} Convierta la lista a un arreglo de Numpy (\texttt{np.array}). Ejecute la misma operación matemática directamente sobre el arreglo completo (vectorización). Mida el tiempo de ejecución.
    \item \textbf{Análisis Crítico:} Calcule cuántas veces más rápido es Numpy en comparación con el bucle tradicional y explique \textbf{por qué} ocurre este fenómeno a nivel de arquitectura (pista: C vs Python y memoria contigua).
\end{enumerate}

\subsection*{Reto 2: Álgebra Lineal Descriptiva (20\%)}
Usted tiene una matriz de $3 \times 3$ que representa las ventas (en miles de USD) de 3 sucursales durante 3 días.
\begin{center}
    $ \text{Ventas} = \begin{bmatrix}
    12 & 15 & 10 \\
    20 & 18 & 22 \\
    8 & 9 & 11
    \end{bmatrix} $
\end{center}
\begin{enumerate}
    \item Utilice \texttt{numpy} para construir la matriz bidimensional (\texttt{np.array}).
    \item Utilice funciones vectorizadas para responder:
    \begin{itemize}
        \item ¿Cuál fue la venta máxima registrada en toda la matriz?
        \item ¿Cuál fue el promedio de ventas por cada \textbf{día} (columnas) (usando el argumento \texttt{axis})?
        \item ¿Cuál fue el total de ventas por \textbf{sucursal} (filas)?
    \end{itemize}
\end{enumerate}

\subsection*{Reto 3: Exploración Sensorial con Pandas (30\%)}
Para este reto, debe investigar cómo generar un DataFrame de juguete usando Pandas.
\begin{enumerate}
    \item Cree un \texttt{pd.DataFrame} de 5 filas y 4 columnas. Las columnas deben ser: \texttt{['id\_sensor', 'fecha', 'temperatura', 'estado']}. Llene los datos con valores inventados realistas (el estado puede ser "Activo" o "Inactivo").
    \item Utilice los métodos \texttt{.head()}, \texttt{.info()} y \texttt{.describe()} sobre el DataFrame e interprete lo que cada uno arroja en una celda de Markdown.
    \item Ejecute la instrucción \texttt{df['temperatura'].mean()} y explique qué tipo de objeto (clase) retorna extraer una sola columna de un DataFrame (Series vs DataFrame).
\end{enumerate}

\subsection*{Reto 4: Filtrado Booleano Inteligente (20\%)}
Del DataFrame creado en el Reto 3, el equipo de mantenimiento necesita una lista de los sensores que presenten fallas.
\begin{enumerate}
    \item Escriba el código en Pandas para filtrar y mostrar \textbf{únicamente} las filas donde el \texttt{estado} sea "Inactivo" \textbf{Y} (\texttt{\&}) la \texttt{temperatura} sea superior a 30 grados.
    \item \textbf{Análisis Crítico:} ¿Por qué en Pandas utilizamos el operador \texttt{\&} en lugar de la palabra reservada \texttt{and} de Python tradicional al hacer filtrado múltiple?
\end{enumerate}

\vspace{1cm}
\begin{center}
    \textit{\textcolor{gray}{`Vectorizar no es solo una optimización de código, es cambiar la mentalidad algorítmica iterativa por una mentalidad algebraica.'}}
\end{center}

\end{document}
